
    %%%%%%%%%%%%%%%%%%%%%%%%%%%%%%%%%%%%%%%%%%%%%%%%%%%%%%%%%%%%%%%%%%%%%%%%%%%%%%%%
    %%% Classe do documento
    \documentclass[12pt, addpoints]{exam}
    \UseRawInputEncoding

    %%%%%%%%%%%%%%%%%%%%%%%%%%%%%%%%%%%%%%%%%%%%%%%%%%%%%%%%%%%%%%%%%%%%%%%%%%%%%%%%
    % preambulo.tex
    %
    % Arquivo de configuração de packages para uso o LaTeX, conforme minhas
    % preferências e modelos pessoais.
    %
    % Para maiores informações, visite:
    %    https://github.com/abrantesasf/latex
    %
    % NÃO ALTERE SE NÃO SOUBER O QUE ESTÁ FAZENDO!
    %%%%%%%%%%%%%%%%%%%%%%%%%%%%%%%%%%%%%%%%%%%%%%%%%%%%%%%%%%%%%%%%%%%%%%%%%%%%%%%%


    %%%%%%%%%%%%%%%%%%%%%%%%%%%%%%%%%%%%%%%%%%%%%%%%%%%%%%%%%%%%%%%%%%%%%%%%%%%%%%%%
    %%% Estruturas de controle:
    \input{static/utils/controles}

    %%%%%%%%%%%%%%%%%%%%%%%%%%%%%%%%%%%%%%%%%%%%%%%%%%%%%%%%%%%%%%%%%%%%%%%%%%%%%%%%
    %%% Compilação condicional em PDF:
    \input{static/utils/pdf}

    %%%%%%%%%%%%%%%%%%%%%%%%%%%%%%%%%%%%%%%%%%%%%%%%%%%%%%%%%%%%%%%%%%%%%%%%%%%%%%%%
    %%% Configurações de layout da página:
    \input{static/utils/layout}

    %%%%%%%%%%%%%%%%%%%%%%%%%%%%%%%%%%%%%%%%%%%%%%%%%%%%%%%%%%%%%%%%%%%%%%%%%%%%%%%%
    %%% Configurações para autores:
    \input{static/utils/autores}

    %%%%%%%%%%%%%%%%%%%%%%%%%%%%%%%%%%%%%%%%%%%%%%%%%%%%%%%%%%%%%%%%%%%%%%%%%%%%%%%%
    %%% Cabeçalho e rodapé:
    %%% Atenção! É MUITO DIFÍCIL ajustar apropriadamente os running headers
    %%% e running footers da primeira vez. Você pode confiar no LaTeX e deixar que
    %%% ele faça o ajuste automaticamente ou pode quebrar a cabeça e
    %%% tentar melhorar algo que o LaTeX já faz muito bem, mesmo que não fique
    %%% exatamente do jeito que você quer. O que você prefere? Se quiser ajustar
    %%% manualmente, descomente a linha abaixo e faça as configurações no arquivo
    %%% "utils/headings.tex".
    %\input{static/utils/headings}

    %%%%%%%%%%%%%%%%%%%%%%%%%%%%%%%%%%%%%%%%%%%%%%%%%%%%%%%%%%%%%%%%%%%%%%%%%%%%%%%%
    %%% Configuração para as fontes do documento:
    %%% Se você quiser usar fontes próprias ou do sistema, precisará ajustar as
    %%% configurações no arquivo "utils/fontes.tex".
    %%% ATENÇÃO: por padrão eu utilizo XeLaTeX com fontes proprietárias projetadas
    %%% por Matthew Butterick, adquiridas em https://mbtype.com, que não podem ser
    %%% distribuídas devido à licença de utilização. Se você usar XeLaTeX, você
    %%% DEVERÁ OBRIGATORIAMENTE ajustar as fontes no arquivo "utils/fontes.tex".
    \input{static/utils/fontes}

    %%%%%%%%%%%%%%%%%%%%%%%%%%%%%%%%%%%%%%%%%%%%%%%%%%%%%%%%%%%%%%%%%%%%%%%%%%%%%%%%
    %%% Sumário:
    \input{static/utils/toc}

    %%%%%%%%%%%%%%%%%%%%%%%%%%%%%%%%%%%%%%%%%%%%%%%%%%%%%%%%%%%%%%%%%%%%%%%%%%%%%%%%
    %%% Matemática:
    \input{static/utils/matematica}

    %%%%%%%%%%%%%%%%%%%%%%%%%%%%%%%%%%%%%%%%%%%%%%%%%%%%%%%%%%%%%%%%%%%%%%%%%%%%%%%%
    %%% Notas de rodapé:
    \input{static/utils/footnotes}

    %%%%%%%%%%%%%%%%%%%%%%%%%%%%%%%%%%%%%%%%%%%%%%%%%%%%%%%%%%%%%%%%%%%%%%%%%%%%%%%%
    %%% Referências cruzadas, links, citações:
    \input{static/utils/crossrefs}

    %%%%%%%%%%%%%%%%%%%%%%%%%%%%%%%%%%%%%%%%%%%%%%%%%%%%%%%%%%%%%%%%%%%%%%%%%%%%%%%%
    %%% Configurações para os hiperlinks em PDF:
    \input{static/utils/pdfconfig}

    %%%%%%%%%%%%%%%%%%%%%%%%%%%%%%%%%%%%%%%%%%%%%%%%%%%%%%%%%%%%%%%%%%%%%%%%%%%%%%%%
    %%% Glossário e índice remissivo:
    \input{static/utils/glossindex}

    %%%%%%%%%%%%%%%%%%%%%%%%%%%%%%%%%%%%%%%%%%%%%%%%%%%%%%%%%%%%%%%%%%%%%%%%%%%%%%%%
    %%% Referências bibliográficas:
    \input{static/utils/bibliografia}

    %%%%%%%%%%%%%%%%%%%%%%%%%%%%%%%%%%%%%%%%%%%%%%%%%%%%%%%%%%%%%%%%%%%%%%%%%%%%%%%%
    %%% Cores:
    \input{static/utils/cores}

    %%%%%%%%%%%%%%%%%%%%%%%%%%%%%%%%%%%%%%%%%%%%%%%%%%%%%%%%%%%%%%%%%%%%%%%%%%%%%%%%
    %%% Computação:
    \input{static/utils/computacao}

    %%%%%%%%%%%%%%%%%%%%%%%%%%%%%%%%%%%%%%%%%%%%%%%%%%%%%%%%%%%%%%%%%%%%%%%%%%%%%%%%
    %%% Gráficos:
    \input{static/utils/graficos}

    %%%%%%%%%%%%%%%%%%%%%%%%%%%%%%%%%%%%%%%%%%%%%%%%%%%%%%%%%%%%%%%%%%%%%%%%%%%%%%%%
    %%% Ícones:
    \input{static/utils/icones}

    %%%%%%%%%%%%%%%%%%%%%%%%%%%%%%%%%%%%%%%%%%%%%%%%%%%%%%%%%%%%%%%%%%%%%%%%%%%%%%%%
    %%% Blocos:
    \input{static/utils/blocos}

    %%%%%%%%%%%%%%%%%%%%%%%%%%%%%%%%%%%%%%%%%%%%%%%%%%%%%%%%%%%%%%%%%%%%%%%%%%%%%%%%
    %%% Boxes coloridos:
    \input{static/utils/colorbox}

    %%%%%%%%%%%%%%%%%%%%%%%%%%%%%%%%%%%%%%%%%%%%%%%%%%%%%%%%%%%%%%%%%%%%%%%%%%%%%%%%
    %%% Tabelas:
    \input{static/utils/tabelas}

    %%%%%%%%%%%%%%%%%%%%%%%%%%%%%%%%%%%%%%%%%%%%%%%%%%%%%%%%%%%%%%%%%%%%%%%%%%%%%%%%
    %%% Ambientes de listas:
    \input{static/utils/listas}

    %%%%%%%%%%%%%%%%%%%%%%%%%%%%%%%%%%%%%%%%%%%%%%%%%%%%%%%%%%%%%%%%%%%%%%%%%%%%%%%%
    %%% Ambientes floats e similares:
    \input{static/utils/floats}

    %%%%%%%%%%%%%%%%%%%%%%%%%%%%%%%%%%%%%%%%%%%%%%%%%%%%%%%%%%%%%%%%%%%%%%%%%%%%%%%%
    %%% Ajustes para a classe EXAM:
    \input{static/utils/exam}

    %%%%%%%%%%%%%%%%%%%%%%%%%%%%%%%%%%%%%%%%%%%%%%%%%%%%%%%%%%%%%%%%%%%%%%%%%%%%%%%%
    %%% Ajustes para textos:
    \input{static/utils/texto}

    %%%%%%%%%%%%%%%%%%%%%%%%%%%%%%%%%%%%%%%%%%%%%%%%%%%%%%%%%%%%%%%%%%%%%%%%%%%%%%%%
    %%% Ambientes, aliases, comandos e customizações em geral para serem
    %%% utilizadas nos documentos. Defina aqui qualquer customização específica
    %%% para seus documentos!
    \input{static/utils/customizacoes}

    %%%%%%%%%%%%%%%%%%%%%%%%%%%%%%%%%%%%%%%%%%%%%%%%%%%%%%%%%%%%%%%%%%%%%%%%%%%%%%%%
    %%% Ajuste do layout, espaçamento de linhas e etc.:
    \geometry{a4paper, portrait, left=2.0cm, right=2.0cm, top=2.0cm, bottom=2.0cm}
    \singlespacing

    %%%%%%%%%%%%%%%%%%%%%%%%%%%%%%%%%%%%%%%%%%%%%%%%%%%%%%%%%%%%%%%%%%%%%%%%%%%%%%%%
    %%% Configurações de cabeçalho e rodapé (não use fancyhdr pois ocorre conflito):
    \pagestyle{headandfoot}
    \runningheadrule
    \firstpageheader{}{}{}
    \runningheader{banco de dados}{}{Avaliação}
    \firstpagefooter{}{}{Página \thepage\ de \numpages}
    \runningfooter{}{}{Página \thepage\ de \numpages}
    \runningfootrule

    %%%%%%%%%%%%%%%%%%%%%%%%%%%%%%%%%%%%%%%%%%%%%%%%%%%%%%%%%%%%%%%%%%%%%%%%%%%%%%%%
    %%% Configurações para as propriedades do PDF:
    \hypersetup{
    hidelinks,           % Comente para web, descomente para impressão
    colorlinks=true,      % True para web, False para impressão
    pdftitle=RicardoMogadoEmBD: Avaliação,
    pdfauthor=Matheus Endlich Silveira,
    pdfsubject=banco de dados,
    pdfkeywords=['Sem tag'],
    pdfinfo={
        CreationDate={}, % Ex.: D:AAAAMMDDHH24MISS
        ModDate={}       % Ex.: D:AAAAMMDDHH24MISS
    }
    }
    \input{static/utils/pdfconfig}

    %%%%%%%%%%%%%%%%%%%%%%%%%%%%%%%%%%%%%%%%%%%%%%%%%%%%%%%%%%%%%%%%%%%%%%%%%%%%%%%%
    %%% Começa o documento
    \begin{document}

    %%%%%%%%%%%%%%%%%%%%%%%%%%%%%%%%%%%%%%%%%%%%%%%%%%%%%%%%%%%%%%%%%%%%%%%%%%%%%%%%
    %%% Folha de instruções padronizadas da UVV para a prova
    \pdfbookmark[1]{Instruções}{instrucoes}
    \begin{figure}[H]
        \begin{center}
            \includegraphics[scale=0.3]{{static/imgs/logo-uvv.png}}
        \end{center}	
    \end{figure}

    \vspace{-1cm}

    \begin{table}[H]
        \centering
        \begin{tabular}{|llll|}
            \hline
            \multicolumn{2}{|l|}{\textbf{Disciplina:} banco de dados} & 
            \multicolumn{1}{l|}{\multirow{3}{*}{\textbf{\makecell{Nota:\\ \,\\ \,}}}} & 
            \multirow{3}{*}{\textbf{\makecell{Coordenador:\\ \,\\ \,}}} \\ 
            \cline{1-2}
            \multicolumn{2}{|l|}{\textbf{Professor:} Matheus Endlich Silveira \hspace{4.72cm} } & 
            \multicolumn{1}{l|}{} & \\ 
            \cline{1-2}
            \multicolumn{2}{|l|}{\textbf{Aluno:}} & 
            \multicolumn{1}{l|}{} & \\ 
            \hline
            \multicolumn{1}{|l|}{\textbf{Turma:} \hspace{6.3cm} } & 
            \multicolumn{1}{l|}{\textbf{Semestre:}} & 
            \multicolumn{2}{l|}{\textbf{Valor:} 10 pontos} \\ 
            \hline
            \multicolumn{1}{|l|}{\textbf{Data:}} & 
            \multicolumn{3}{l|}{\textbf{Avaliação:}} \\ 
            \hline
        \end{tabular}
    \end{table}

    \begin{center}
        \fbox{\setlength\fboxsep{0.3cm} \fbox{\parbox{15.4cm}{
            \begin{center}
                \textbf{INSTRUÇÕES PARA A REALIZAÇÃO DESTA AVALIAÇÃO:}
            \end{center}
            \begin{itemize}[leftmargin=*]
                \item Esta é a primeira avaliação da disciplina banco de dados, 
                    e engloba o conteúdo referente Sem tag.
                \item Esta avaliação contém 20 questões objetivas, \textbf{todas obrigatórias}.
                \item \textbf{Leia atentamente cada questão!} As questões objetivas terão uma,
                    e apenas uma única, resposta a ser marcada.
                \item \textbf{Desligue o celular} antes de começar. A avaliação é
                    \textbf{individual} e \textbf{sem consulta}.
                \item As questões podem ser respondidas, na prova, com lápis ou caneta. Ao final
                    da prova \textbf{{VOCÊ DEVERÁ PREENCHER O GABARITO COM CANETA
                    DE COR PRETA}} \textbf{\underline{OU AZUL}} para que seja feita a correção
                    automatizada através da leitura do padrão de respostas marcado no
                    gabarito.
                \item \textbf{CUIDADO ao preencher o gabarito} pois eles são \textbf{INDIVIDUAIS
                    e NOMEADOS} (cada estudante tem um gabarito próprio específico, com um
                    código de identificação único). \textbf{Questões rasuradas são anuladas}
                    pelo software de leitura do gabarito.
                \item Ao final da prova, \textbf{DEVOLVA A PROVA E O GABARITO} para o professor.
                \item Siga todas as normas de \textbf{Integridade Acadêmica} da disciplina.
                    Alunos que forem flagrados com qualquer espécie de ``cola'' ou trocando
                    informações com outros alunos terão suas avaliações recolhidas, as notas
                    zeradas, e a situação será encaminhada para a coordenação para a aplicação
                    das penalidades previstas pela universidade.
                \item Com 90 minutos de avaliação você terá tempo de até 4,min por questão,
                    em média.
                \item Boa avaliação!
            \end{itemize}
            \vspace{2cm}
        }}}
    \end{center}
    \newpage

    %%%%%%%%%%%%%%%%%%%%%%%%%%%%%%%%%%%%%%%%%%%%%%%%%%%%%%%%%%%%%%%%%%%%%%%%%%%%%%%%
    %%% Inicia o bloco de questões:
    \begin{questions}

    \newpage
    \question

Qual das seguintes opções \textbf{não é} uma vantagem do uso de um banco de

dados sobre o sistema de arquivos convencional?


\begin{choices}
\choice Controle de redundância e inconsistência
\choice Independência de dados
\choice Eficiência
\choice Controle de acesso e visualização de dados
\choice Estruturação de dados
\end{choices}

\question

Por que não podemos utilizar apelidos de coluna na cláusula WHERE da SQL?
\begin{choices}
\choice Porque a cláusula WHERE é avaliada antes da SELECT.
\choice Porque só podemos utilizar aliases, e não apelidos.
\choice Porque a cláusula SELECT é avaliada antes da WHERE.
\choice Porque os apelidos só são avaliados depois da cláusula ORDER BY.
\choice Nós podemos utilizar apelidos na cláusula WHERE.
\end{choices}

\question

Analise o diagrama relacional do banco de dados exibido na figura abaixo:

\begin{figure}[H]

  \centering

  \caption{Banco de dados de peças e fornecedores}

  \label{fig:pecas}

  %\fbox{

    \includegraphics[width=0.5\linewidth]{static/imgs/defaul.png}

  %}\\

  %\footnotesize{Fonte: xxxx.}

\end{figure}

Podemos dizer que esse banco de dados ilustra:
\begin{choices}
\choice Um relacionamento com cardinalidade N:N entre as tabelas ``pecas' e ``fornecedores', realizado através de relacionamentos identificados através da tabela ``fornecimentos'. 
\choice Um relacionamento com identificação N:N entre as tabelas ``pecas' e ``fornecedores', realizado através do relacionamento com cardinalidade identificada através da tabela ``fornecimentos'. 
\choice Um relacionamento com cardinalidade 1:N entre as tabelas ``pecas' e ``fornecedores', indicando que uma peça pode ter sido fornecida por diversos fornecedores. 
\choice Um relacionamento com cardinalidade N:N entre as tabelas ``pecas' e ``fornecimentos', realizado através de relacionamentos não identificados com a tabela ``fornecedores'. 
\choice Um relacionamento com cardinalidade N:N entre as tabelas ``pecas' e ``fornecedores', realizado através de relacionamentos não identificados através da tabela ``fornecimentos'. 
\end{choices}

\question

São modelos de dados legados, pré-relacionais, exceto:


\begin{choices}
\choice Objeto
\choice Nenhuma das alternativas anteriores
\choice Flat file
\choice Rede
\choice Hierárquico
\end{choices}

\question

Dentre as definições a seguir, ligadas ao conceito de normalização do modelo relacional, qual delas é \textbf{INCORRETA}?


\begin{choices}
\choice A normalização é um processo passo a passo reversível de substituição de uma dada coleção de relações por sucessivas coleções de relações as quais possuem uma estrutura progressivamente mais simples e mais regular.
\choice As relações que obedecem à primeira forma normal não apresentam anomalias.
\choice A primeira forma normal estabelece que os atributos da relação contêm apenas valores atômicos.
\choice O objetivo da normalização é eliminar várias anomalias (ou aspectos indesejáveis) de uma relação.
\choice As formas normais se baseiam em certas estruturas de dependências.
\end{choices}

\question

(ENADE 2017) No modelo relacional de dados, uma junção entre tabelas cria uma

outra tabela derivada de duas ou mais tabelas de acordo com os critérios

especificados para a junção, de modo similar às regras da teoria dos

conjuntos. Nos conjuntos a seguir, as áreas hachuradas representam os resultados

das consultas SQL em que ``id' é uma chave primária e ``fk' é uma chave

estrangeira:

\begin{figure}[H]

  \centering

  \includegraphics[width=0.5\linewidth]{static/imgs/defaul.png}

\end{figure}

Assinale a opção em que a declaração SQL é equivalente ao conjunto apresentado a

seguir:

\begin{figure}[H]

  \centering

  \includegraphics[width=0.5\linewidth]{static/imgs/defaul.png}

\end{figure}


\begin{choices}
\choice \begin{verbatim}SELECT *

FROM tabela1 t1

RIGHT JOIN tabela2 t2 ON (t1.id = t2.fk)

WHERE t1.id IS NULL;

\end{verbatim}


\choice \begin{verbatim}SELECT *

FROM tabela1 t1

RIGHT JOIN tabela2 t2 ON (t1.id = t2.fk)

WHERE t2.fk <> NULL;

\end{verbatim}


\choice \begin{verbatim}SELECT *

FROM tabela1 t1

WHERE NOT EXISTS (SELECT 1

                  FROM tabela1 t2

                  WHERE t2.id = tk.fk);

\end{verbatim}


\choice \begin{verbatim}SELECT *

FROM tabela1 t1

INNER JOIN tabela2 t2 ON (t1.id = t2.fk);

\end{verbatim}


\choice \begin{verbatim}SELECT *

FROM tabela1 t1 WHERE EXISTS (SELECT 1

                              FROM tabela2 t2

                              WHERE t2.id = t1.fk);

\end{verbatim}


\end{choices}

\question

Considere as seguintes tabelas em uma base de dados relacional (chaves primárias sublinhadas):\\

Departamento (\underline{CodDepto}, NomeDepto)\\

Empregado (\underline{CodEmp}, NomeEmp, CodDepto)\\

Considere as seguintes restrições de integridade sobre esta base de dados relacional:\\

- Empregado.CodDepto é sempre diferente de NULL\\

- Empregado.CodDepto é chave estrangeira da tabela Departamento com cláusulas ON DELETE RESTRICT e ON UPDATE RESTRICT\\

Qual das seguintes validações não é especificada por estas restrições de integridade:
\begin{choices}
\choice Sempre que uma linha for excluída de Departamento, deve ser garantido que o valor de Departamento.CodDepto não aparece na coluna Empregado.CodDepto.
\choice Sempre que o valor de Empregado.CodDepto for alterado, deve ser garantido que o novo valor de Empregado.CodDepto aparece em Departamento.CodDepto.
\choice Sempre que uma nova linha for inserida em Empregado, deve ser garantido que o valor de Empregado.CodDepto aparece na coluna Departamento.CodDepto.
\choice Sempre que o valor de Departamento.CodDepto for alterado, deve ser garantido que não há uma linha com o antigo valor de Departamento.CodDepto na coluna Empregado.CodDepto.
\choice Sempre que uma nova linha for inserida em Departamento, deve ser garantido que o valor de Departamento.CodDepto aparece na coluna Empregado.CodDepto.
\end{choices}

\question

Ao analisar um novo sistema de gerenciamento de bancos de dados que surgiu no

mercado, você percebe que a entrada às operações do modelo de dados são tabelas

mas as saídas nunca são tabelas, são sempre outras estruturas. O vendedor desse

SGBD afirma que ele é um SGBD relacional. Pode-se afirmar que:


\begin{choices}
\choice Esse SGBD é relacional pois consegue trabalhar com tabelas para a entrada das operações do modelo de dados. A saída não é importante.
\choice Nenhuma das alternativas anteriores.
\choice Esse SGBD é relacional pois o vendedor afirmou que é.
\choice Esse SGBD não é relacional pois, para ser verdadeiramente relacional, as operações do modelo de dados não podem trabalhar recebendo tabelas como entrada.
\choice Esse SGBD não é relacional pois no modelo relacional a única estrutura percebida pelos usuários são tabelas.
\end{choices}

\question

Um banco de dados (BD) é uma coleção logicamente coerente de dados relacionados

e persistentes. Em relação às características de um banco de dados, assinale a

alternativa verdadeira:
\begin{choices}
\choice Não precisam ser obrigatoriamente computadorizados, existem bancos de dados manuais e/ou em papel.
\choice Geralmente é projetado e usado para finalidades gerais, sendo extremamente genéricos e adaptáveis para qualquer situação.
\choice Armazena dois grandes ``tipos`` de dados: os dados do usuário final e os metadados, sendo que os metadados ficam armazenados nas mesmas tabelas que os dados dos usuários e isso nos permite saber a estrutura dos bancos de dados.
\choice Representa um aspecto (ou parte) do mundo real, parte essa chamada de ``catálogo' ou ``dicionário de dados``.
\choice Não precisa de dados atuais e verdadeiros para ser útil (para representar corretamente a realidade atual).
\end{choices}

\question

(Adaptado do ENADE 2005) Avalie o seguite diagrama relacional de um banco de dados:

\begin{figure}[H]

  \centering

  \caption{Banco de dados de peças e fornecedores}

  \label{fig:pecas}

  %\fbox{

    \includegraphics[width=0.5\linewidth]{static/imgs/defaul.png}

  %}\\

  %\footnotesize{Fonte: xxxx.}

\end{figure}

Assinale a opção que apresenta um comando SQL que permite obter uma lista que

contenha o nome de cada fornecedor que tenha fornecido alguma peça, o código da

peça fornecida, a descrição dessa peça e a quantidade fornecida da referida

peça.
\begin{choices}
\choice \begin{verbatim}SELECT DISTINCT p.nome, p.codigo, p.descricao, f.quantidade FROM pecas p INNER JOIN fornecimentos f ON (p.codigo = f.cod_peca);\end{verbatim} 
\choice \begin{verbatim}SELECT * FROM pecas INNER JOIN fornecimentos f ON (p.codigo = f.cod_peca) INNER JOIN fornecedores f2 ON (f.cod_forn = f2.cod_forn) ORDER BY f2.nome;\end{verbatim} 
\choice \begin{verbatim}SELECT nome, codigo, descricao, quantidade FROM pecas INNER JOIN FORNECEDORES INNER JOIN FORNECIMENTOS;\end{verbatim} 
\choice \begin{verbatim}SELECT * FROM pecas INNER JOIN fornecedores INNER JOIN fornecimentos;\end{verbatim} 
\choice \begin{verbatim}SELECT f2.nome, p.codigo, p.descricao, f.quantidade FROM pecas INNER JOIN fornecimentos f ON (p.codigo = f.cod_peca) INNER JOIN fornecedores f2 ON (f.cod_forn = f2.cod_forn);\end{verbatim} 
\end{choices}

\question

Na linguagem SQL, qual das seguintes palavras-chave é usada para modificar os

dados em uma tabela existente?


\begin{choices}
\choice Nenhuma das alternativas acima
\choice UPDATE
\choice INSERT
\choice SELECT
\choice CREATE
\end{choices}

\question

(ENADE 2014) Considere o seguinte banco de dados ilustrado no diagrama relaconal

abaixo:

\begin{figure}[H]

  \centering

  \caption{Estados e fornecedores}

  \label{fig:estados}

  %\fbox{

    \includegraphics[width=0.5\linewidth]{static/imgs/defaul.png}

  %}\\

  %\footnotesize{Fonte: xxxx.}

\end{figure}

A expressão SQL que obtém os nomes dos estados para os quais não há fornecedores

cadastrados é:


\begin{choices}
\choice \begin{verbatim}SELECT e.nome

FROM estado e,

FROM fornecedor f

WHERE e.uf = f.uf;

\end{verbatim}


\choice \begin{verbatim}SELECT e.nome

FROM estado e

WHERE e.uf IN (SELECT f.uf

               FROM fornecedor f);

\end{verbatim}


\choice \begin{verbatim}SELECT e.nome

FROM estado e

WHERE e.uf NOT IN (SELECT f.uf

                   FROM fornecedor f);

\end{verbatim}


\choice \begin{verbatim}SELECT e.nome

FROM estado e,

FROM fornecedor f

WHERE e.nome = f.uf;

\end{verbatim}


\choice \begin{verbatim}SELECT e.uf

FROM estado e

WHERE e.nome NOT IN (SELECT f.uf

                     FROM fornecedor f);

\end{verbatim}
\end{choices}

\question

Seu chefe que fazer uma campanha de venda para produtos que custam entre 20 e 30

reais, mas está preocupado porque ouviu de um dos vendedores que diversos

produtos nessa faixa de preço estão com o estoque zerado. Para verificar a

situação seu chefe pediu para que você prepare um relatório que mostre os

produtos que custam entre 20 e 30 reais e que estão com o estoque zerado. Você

preparou o seguinte relatório para seu chefe:

\begin{tcolorbox}

 \begin{verbatim}1  SELECT l.nome AS loja

2       , p.nome AS produto

3       , p.preco_unitario AS preco

4       , e.quantidade AS qtd

5  FROM produtos p

6  INNER JOIN estoques e ON (e.produto_id = p.produto_id)

7  INNER JOIN lojas l ON (e.loja_id = l.loja_id)

8  WHERE (p.preco_unitario >= 20 OR p.preco_unitario <= 30)

9    AND e.quantidade = 0

10 ORDER BY loja, produto;

 \end{verbatim}

\end{tcolorbox}

O seu relatório:
\begin{choices}
\choice Está errado pois há um erro de sintaxe nas linhas 6 e 7.
\choice Está correto e trará exatamente, os produtos com o estoque zerado e na faixa de preço requisitada pelo seu chefe. Além disso ordena o resultado pelo nome da loja e pelo nome do produto.
\choice Está errado pois a tabela que deveria estar na cláusula FROM é a tabela ``lojas', não a tabela ``produtos'.
\choice Está errado pois vai trazer também produtos fora da faixa de preço que o seu chefe quer.
\choice Nenhuma das alternativas anteriores.
\end{choices}

\question

Considere a seguinte tabela em uma base de dados relacional (chave primária sublinhada):\\

Tabela1 (CodAluno, CodDisciplina, AnoSemestre, NomeAluno, NomeDisciplina, CodNota, DescricaoNota)\\

Considere as seguintes dependências funcionais:

\begin{itemize}

    \item CodAluno \(\rightarrow\) NomeAluno

    \item CodDisciplina \(\rightarrow\) NomeDisciplina

    \item (CodAluno, CodDisciplina, AnoSemestre) \(\rightarrow\) CodNota

    \item (CodAluno, CodDisciplina, AnoSemestre) \(\rightarrow\) DescricaoNota

    \item CodNota \(\rightarrow\) DescricaoNota

\end{itemize}



Considerando as formas normais, qual das afirmativas abaixo se aplica:
\begin{choices}
\choice A tabela encontra-se na terceira forma normal, mas não na quarta forma normal.
\choice A tabela está na quarta forma normal.
\choice A tabela não está na primeira forma normal.
\choice A tabela encontra-se na segunda forma normal, mas não na terceira forma normal.
\choice A tabela encontra-se na primeira forma normal, mas não na segunda forma normal.
\end{choices}

\question

``Um modelo de dados é uma definição lógica, abstrata e auto-contida

dos \circled{1}, \circled{2} e \circled{3} que, juntos, constituem a \circled{4}

com a qual os usuários interagem. Os \circled{1} nos permitem modelar

a \circled{5}; os \circled{2} nos permitem modelar

seu \circled{6}. Os \circled{3} nos permitem modelar sua \circled{7}.'



As palavras/termos que completam a frase acima, são, conforme os números:


\begin{choices}
\choice 1 = operadores; 2 = objetos; 3 = demais conceitos; 4 = máquina abstrata; 5 = estrutura de dados; 6 = comportamento; 7 = integridade.
\choice 1 = objetos; 2 = operadores; 3 = máquina abstrata; 4 = demais conceitos; 5 = estrutura de dados; 6 = integridade; 7 = comportamento.
\choice 1 = objetos; 2 = máquina abstrata; 3 = demais conceitos; 4 = operadores; 5 = estrutura de dados; 6 = comportamento; 7 = integridade.
\choice 1 = demais conceitos; 2 = objetos; 3 = operadores; 4 = estrutura de dados; 5 = máquina abstrata; 6 = integridade; 7 = comportamento.
\choice 1 = objetos; 2 = operadores; 3 = demais conceitos; 4 = máquina abstrata; 5 = estrutura de dados; 6 = comportamento; 7 = integridade.
\end{choices}

\question

Um banco de dados OLAP (Online Analytical Processing --- Processamento Analítico

Online) é um banco de dados organizado de forma a dar suporte ao Business

Intelligence (BI --- Inteligência de Negócio) que, de forma simplificada, é o

processo de analisar dados para obter informações que podem ser utilizadas na

tomada de decisões comerciais e estratégicas. Em relação aos bancos de dados

OLAP, marque a afirmação correta:
\begin{choices}
\choice São preparados para responder perguntas do tipo ``Quais alunos estão matriculados na turma de Bancos de Dados I neste semestre?'
\choice São otimizados para consulta e relatórios de rotina de dados.
\choice São otimizados para o processamento de grande volume de dados históricos e métricas de negócio para a tomada de decisões estratégicas, permitindo e facilitando a análise de dados para a geração de informações e conhecimentos para a tomada de decisões.
\choice Geralmente os dados de origem do OLTP são de bancos de dados OLAP.
\choice São otimizados para o processamento de grande volume de transações diárias, o que permite a análise e a tomada de decisões em tempo real, favorecendo assim a administração empresarial.
\end{choices}

\question

O usuário que não é da área de tecnologia, costuma ser de alta gerência, que é

especialista no negócio e tem a responsabilidade central pelas políticas de

dados da empresa é o:
\begin{choices}
\choice Encarregado da Proteção de Dados (DPO: Data Protection Officer)
\choice Coordenador de Tecnologia da Informação (ITS: IT Supervisor)
\choice Gerente de Tecnologia da Informação (ITM: IT Manager)
\choice Administrador de Dados (DA: Data Administrator)
\choice Administrador do Banco de Dados (DBA: Database Administrator)
\end{choices}

\question

(Adaptado do ENADE 2005) Analise o relacionamento existente entre as tabelas

``clubes' e ``jogadores', nas alternativas abaixo, e marque a alternativa que

indica que todo jogador deve pertencer a um único clube, e todo clube poder ter

ou não jogadores.
\begin{choices}
\choice \includegraphics[width=0.5\linewidth]{static/imgs/defaul.png} \vspace{0.5cm}
\choice \includegraphics[width=0.5\linewidth]{static/imgs/defaul.png} \vspace{0.5cm}
\choice \includegraphics[width=0.5\linewidth]{static/imgs/defaul.png} \vspace{0.5cm}
\choice \includegraphics[width=0.5\linewidth]{static/imgs/defaul.png} \vspace{0.5cm}
\choice \includegraphics[width=0.5\linewidth]{static/imgs/defaul.png} \vspace{0.5cm}
\end{choices}

\question

Você estava ouvindo a discussão de dois colegas, o João e a Maria. Enquanto o

João afirmava que o conceito de banco de dados (BD) é extremamente diferente do

conceito de sistema de gerenciamento de bancos de dados (SGBD), a Maria

discordava e afirmava que o banco de dados nada mais é do que um componente do

SGBD. Nesse momento o professor entrou na sala, percebeu a discussão e falou

que:
\begin{choices}
\choice O João estava certo já que o banco de dados é, realmente, um conceito diferente do conceito do sistema de gerenciamento de bancos de dados, e existe de forma independente.
\choice Os dois estavam certos pois o conceito, na prática, depende da situação específica que se está considerando.
\choice Não é possível saber quem estava certo ou errado pois não definiram qual o modelo de dados que estava sendo considerado na discusão.
\choice Os dois estavam errados pois não são conceitos separados nem componentes, são conceitos sinônimos.
\choice A Maria estava certa já que o banco de dados é, realmente, um componente interno do sistema de gerenciamento de bancos de dados e não existe de forma independente.
\end{choices}

\question

Dentre as definições a seguir, ligadas ao conceito de visões do modelo relacional, qual delas é \textbf{INCORRETA}?


\begin{choices}
\choice Uma visão relacional é uma relação virtual, derivada de relações base a partir da especificação de operações da álgebra relacional.
\choice Uma visão relacional é uma relação virtual que nunca é materializada.
\choice O gerenciamento de visões envolve a conversão da consulta do usuário sobre as visões para a consulta sobre as relações base.
\choice Programas aplicativos do banco de dados podem ser executados sobre visões de relações da base de dados.
\choice Uma visão é útil por representar uma percepção particular do banco de dados, compartilhado por muitos aplicativos.
\end{choices}



    %%%%%%%%%%%%%%%%%%%%%%%%%%%%%%%%%%%%%%%%%%%%%%%%%%%%%%%%%%%%%%%%%%%%%%%%%%%%%%%%
    %%% Finaliza o bloco de questões:
    \end{questions}
    \end{document}
    